% --- CORRECCIÓN IMPORTANTE EN LA PRIMERA LÍNEA ---
% Agregamos "xcolor=table" para que funcione \rowcolor sin errores
\documentclass[aspectratio=169, 10pt, xcolor=table]{beamer}

% --- Configuración del Tema (estilo profesional como el ejemplo) ---
\usetheme{Madrid}
\usecolortheme{dolphin}

% --- Paquetes necesarios ---
\usepackage[utf8]{inputenc}
\usepackage[spanish]{babel}
\usepackage{amsmath, amssymb}
\usepackage{booktabs}
\usepackage{graphicx}
\usepackage{listings}
\usepackage{tikz}
\usetikzlibrary{arrows.meta, positioning, shapes.geometric}

% --- Ruta de Imágenes (crea la carpeta "imagenes" y coloca ahí tus archivos) ---
\graphicspath{{./imagenes/}}

% --- Estilo para código Python ---
\definecolor{codegreen}{rgb}{0,0.6,0}
\definecolor{codegray}{rgb}{0.5,0.5,0.5}
\definecolor{codepurple}{rgb}{0.58,0,0.82}
\definecolor{backcolour}{rgb}{0.95,0.95,0.92}
\lstset{
    language=Python,
    backgroundcolor=\color{backcolour},
    commentstyle=\color{codegreen},
    keywordstyle=\color{blue},
    numberstyle=\tiny\color{codegray},
    stringstyle=\color{codepurple},
    basicstyle=\ttfamily\scriptsize,
    breaklines=true,
    frame=single,
    showstringspaces=false
}

% --- Pie de página personalizado ---
\makeatletter

\setbeamertemplate{footline}{
  \leavevmode%
  \hbox{%
  \begin{beamercolorbox}[wd=.333333\paperwidth,ht=2.25ex,dp=1ex,center]{author in head/foot}%
    \usebeamerfont{author in head/foot}\insertshortauthor
  \end{beamercolorbox}%
  \begin{beamercolorbox}[wd=.333333\paperwidth,ht=2.25ex,dp=1ex,center]{title in head/foot}%
    \usebeamerfont{title in head/foot}Sesión 4
  \end{beamercolorbox}%
  \begin{beamercolorbox}[wd=.333333\paperwidth,ht=2.25ex,dp=1ex,right]{date in head/foot}%
    \usebeamerfont{date in head/foot}Machine Learning \hfill \insertframenumber/\inserttotalframenumber \hspace*{0.3cm}
  \end{beamercolorbox}}%
  \vskip0pt%
}

\makeatother

% --- Datos de la portada ---
\title[SESIÓN 4]{\textbf{Sesión 4}}
\subtitle{Clasificación I: Regresión Logística y Métricas}
\author[Prof. C. Sánchez]{Carlos César Sánchez Coronel}
\institute{\textbf{MACHINE LEARNING}}
\date{}

% Logo opcional (descomentar si tienes la imagen)
% \titlegraphic{\includegraphics[height=1.5cm]{logo.png}}

% --- Eliminar la barra de navegación (opcional) ---
\setbeamertemplate{navigation symbols}{}

% --- Índice al inicio de cada sección ---
\AtBeginSection[]{
  \begin{frame}
    \frametitle{Contenido}
    \tableofcontents[currentsection]
  \end{frame}
}

% ============================================================================
% INICIO DEL DOCUMENTO
% ============================================================================
\begin{document}

% --- Portada ---
\begin{frame}
    \titlepage
\end{frame}

% --- Índice general ---
\begin{frame}{Contenido}
    \tableofcontents
\end{frame}

% ============================================================================
% SECCIÓN 1: INTRODUCCIÓN A LA CLASIFICACIÓN
% ============================================================================
\section{Introducción a la Clasificación}

\begin{frame}{¿Por qué no regresión lineal para clasificación?}
    \begin{itemize}
        \item Predicciones fuera del rango $[0,1]$ (no interpretables como probabilidad).
        \item Errores no homocedásticos (varianza variable).
        \item Relación no lineal; la sigmoide es más adecuada.
    \end{itemize}
    \vspace{0.3cm}
    \begin{exampleblock}{Conclusión}
        La regresión logística modela directamente la probabilidad de pertenencia a una clase.
    \end{exampleblock}
    % Basado en Pregunta 1
    % IMAGEN: Comparación de regresión lineal vs logística
\end{frame}

% ============================================================================
% SECCIÓN 2: REGRESIÓN LOGÍSTICA
% ============================================================================
\section{Regresión Logística}

\begin{frame}{Función Sigmoide}
    \begin{block}{Definición}
        \[
        \sigma(z) = \frac{1}{1 + e^{-z}}, \quad z \in \mathbb{R}
        \]
        \begin{itemize}
            \item $\sigma(z) \to 0$ cuando $z \to -\infty$
            \item $\sigma(z) \to 1$ cuando $z \to +\infty$
            \item $\sigma(0) = 0.5$
            \item Monótona creciente y diferenciable.
        \end{itemize}
    \end{block}
    % Basado en Pregunta 2
    % IMAGEN: Gráfica de la sigmoide
\end{frame}

\begin{frame}{Modelo Probabilístico}
    \begin{block}{Probabilidad de clase positiva}
        \[
        P(y=1|\mathbf{x}) = \sigma(\beta_0 + \beta_1 x_1 + \dots + \beta_p x_p) = \frac{1}{1 + e^{-(\beta_0 + \boldsymbol{\beta}^T \mathbf{x})}}
        \]
        $P(y=0|\mathbf{x}) = 1 - P(y=1|\mathbf{x})$.
    \end{block}
    \begin{exampleblock}{Interpretación con log-odds}
        \[
        \text{logit}(p) = \ln\left(\frac{p}{1-p}\right) = \beta_0 + \boldsymbol{\beta}^T \mathbf{x}
        \]
        Cada $\beta_j$ mide el cambio en log-odds por unidad de $x_j$.
    \end{exampleblock}
    % Basado en Pregunta 3
\end{frame}

\begin{frame}{Interpretación de Coeficientes: Odds Ratio}
    \begin{itemize}
        \item \textbf{Odds ratio:} $e^{\beta_j}$.
        \item Si $e^{\beta_j} > 1$, la variable aumenta las odds de la clase positiva.
        \item Si $e^{\beta_j} < 1$, la variable disminuye las odds.
    \end{itemize}
    \begin{exampleblock}{Ejemplo}
        En un modelo de riesgo crediticio, $\beta_{\text{ingresos}} = -0.05$ implica $e^{-0.05} \approx 0.951$, es decir, por cada unidad de ingreso adicional, las odds de default se reducen en un $4.9\%$.
    \end{exampleblock}
    % Basado en Pregunta 3
\end{frame}

\begin{frame}{Función de Pérdida: Log-Loss (Entropía Cruzada)}
    \begin{block}{Para una muestra}
        \[
        \mathcal{L}(\hat{p}, y) = -[y \log(\hat{p}) + (1-y) \log(1-\hat{p})]
        \]
    \end{block}
    \begin{block}{Para todo el conjunto}
        \[
        J(\boldsymbol{\beta}) = -\frac{1}{n} \sum_{i=1}^n \left[ y_i \log(\hat{p}_i) + (1-y_i) \log(1-\hat{p}_i) \right]
        \]
        \begin{itemize}
            \item Convexa → mínimo global único.
            \item Diferenciable → apta para gradiente descendente.
        \end{itemize}
    \end{block}
    % Basado en Pregunta 4
\end{frame}

\begin{frame}{Optimización: Gradiente Descendente}
    \begin{block}{Gradiente de la log-loss}
        \[
        \nabla J(\boldsymbol{\beta}) = \frac{1}{n} \sum_{i=1}^n (\hat{p}_i - y_i) \mathbf{x}_i
        \]
        Misma forma que en regresión lineal, pero con $\hat{p}_i$ definido por la sigmoide.
    \end{block}
    \begin{itemize}
        \item Actualización: $\boldsymbol{\beta} := \boldsymbol{\beta} - \alpha \nabla J(\boldsymbol{\beta})$.
        \item Se pueden usar variantes: SGD, mini-batch.
    \end{itemize}
    % Basado en Pregunta 5
\end{frame}

\begin{frame}{Límite de Decisión}
    \begin{block}{Definición}
        Puntos donde $P(y=1|\mathbf{x}) = 0.5$, es decir:
        \[
        \beta_0 + \boldsymbol{\beta}^T \mathbf{x} = 0
        \]
        \begin{itemize}
            \item En 2D: una recta.
            \item En 3D: un plano.
            \item En dimensiones superiores: un hiperplano.
        \end{itemize}
    \end{block}
    \begin{exampleblock}
        La regresión logística es un \textbf{clasificador lineal}.
    \end{exampleblock}
    % Basado en Pregunta 2
    % IMAGEN: Límite de decisión lineal en 2D
\end{frame}

\begin{frame}{Regularización en Regresión Logística}
    \begin{columns}[t]
        \column{0.5\textwidth}
        \begin{block}{Ridge (L2)}
            $J_{\text{ridge}} = J + \lambda \sum \beta_j^2$
            \begin{itemize}
                \item Encoge coeficientes, pero no los anula.
            \end{itemize}
        \end{block}
        \column{0.5\textwidth}
        \begin{block}{Lasso (L1)}
            $J_{\text{lasso}} = J + \lambda \sum |\beta_j|$
            \begin{itemize}
                \item Puede fijar coeficientes a cero → selección de variables.
            \end{itemize}
        \end{block}
    \end{columns}
    \vspace{0.3cm}
    \begin{itemize}
        \item Útil con muchas variables para evitar sobreajuste.
        \item Elastic Net combina ambas.
    \end{itemize}
    % Basado en Pregunta 6
\end{frame}

% ============================================================================
% SECCIÓN 3: EVALUACIÓN DE CLASIFICADORES
% ============================================================================
\section{Evaluación de Clasificadores Binarios}

\begin{frame}{Matriz de Confusión}
    \begin{table}
        \centering
        \begin{tabular}{c|c|c}
            & \textbf{Real Positivo} & \textbf{Real Negativo} \\
            \hline
            \textbf{Pred. Positivo} & VP (Verdaderos Positivos) & FP (Falsos Positivos) \\
            \textbf{Pred. Negativo} & FN (Falsos Negativos) & VN (Verdaderos Negativos) \\
        \end{tabular}
    \end{table}
    \begin{itemize}
        \item VP: spam correcto; FP: legítimo marcado como spam.
        \item FN: spam no detectado; VN: legítimo correcto.
    \end{itemize}
    % Basado en Pregunta 7
    % IMAGEN: Matriz de confusión con colores
\end{frame}

\begin{frame}{Métricas Derivadas}
    \begin{columns}[t]
        \column{0.5\textwidth}
        \begin{block}{Exactitud (Accuracy)}
            $\frac{VP+VN}{VP+VN+FP+FN}$
        \end{block}
        \begin{block}{Precisión (Precision)}
            $\frac{VP}{VP+FP}$
        \end{block}
        \column{0.5\textwidth}
        \begin{block}{Recall (Sensibilidad)}
            $\frac{VP}{VP+FN}$
        \end{block}
        \begin{block}{F1-Score}
            $2 \cdot \frac{\text{Prec} \cdot \text{Recall}}{\text{Prec} + \text{Recall}}$
        \end{block}
    \end{columns}
    \vspace{0.3cm}
    \begin{exampleblock}{Ejemplo fraude (Pregunta 8)}
        VP=15, FP=5, FN=5, VN=75 → Precision=0.75, Recall=0.75, F1=0.75.
    \end{exampleblock}
\end{frame}

\begin{frame}{Comparación de Métricas (Ejemplo Médico)}
    Dos modelos para enfermedad rara (prevalencia 1\%):
    \begin{table}
        \centering
        \begin{tabular}{l|c|c|c}
            \textbf{Modelo} & \textbf{Precisión} & \textbf{Recall} & \textbf{F1} \\
            \hline
            A & 0.444 & 0.80 & 0.57 \\
            B & 0.091 & 1.00 & 0.17 \\
        \end{tabular}
    \end{table}
    \begin{itemize}
        \item Si FN es muy costoso → elegir B (recall 1.0).
        \item Si FP es muy costoso → elegir A (mejor precisión).
    \end{itemize}
    % Basado en Pregunta 9
\end{frame}

\begin{frame}{Curva ROC y AUC}
    \begin{columns}[t]
        \column{0.5\textwidth}
        \begin{block}{Curva ROC}
            \begin{itemize}
                \item Eje X: Tasa de Falsos Positivos (FPR)
                \item Eje Y: Tasa de Verdaderos Positivos (TPR = Recall)
                \item Cada punto = un umbral diferente.
            \end{itemize}
        \end{block}
        \column{0.5\textwidth}
        \begin{block}{AUC (Area Under Curve)}
            \begin{itemize}
                \item AUC = 0.5: clasificador aleatorio.
                \item AUC = 1.0: clasificador perfecto.
                \item Probabilidad de que un positivo tenga score mayor que un negativo.
            \end{itemize}
        \end{block}
    \end{columns}
    % Basado en Pregunta 10
    % IMAGEN: Curva ROC con AUC
\end{frame}

\begin{frame}{Cuándo usar cada métrica}
    \begin{itemize}
        \item \textbf{Accuracy:} clases balanceadas, costos simétricos.
        \item \textbf{Precision:} minimizar falsos positivos (ej. spam, diagnósticos graves).
        \item \textbf{Recall:} minimizar falsos negativos (ej. fraude, screening).
        \item \textbf{F1:} balance entre precisión y recall, clases desbalanceadas.
        \item \textbf{AUC-ROC:} comparación global de modelos independiente del umbral.
    \end{itemize}
    \vspace{0.3cm}
    \begin{alertblock}{¡Cuidado!}
        En gran desbalance, AUC-ROC puede ser optimista; usar curva Precision-Recall (PR-AUC).
    \end{alertblock}
    % Basado en Pregunta 11
\end{frame}

% ============================================================================
% SECCIÓN 4: CLASES DESBALANCEADAS
% ============================================================================
\section{Manejo de Clases Desbalanceadas}

\begin{frame}{Técnicas de Muestreo}
    \begin{columns}[t]
        \column{0.33\textwidth}
        \begin{block}{Undersampling}
            Reducir clase mayoritaria.
            \begin{itemize}
                \item Ventaja: reduce datos.
                \item Desventaja: pierde información.
            \end{itemize}
        \end{block}
        \column{0.33\textwidth}
        \begin{block}{Oversampling}
            Duplicar clase minoritaria.
            \begin{itemize}
                \item Puede causar sobreajuste.
            \end{itemize}
        \end{block}
        \column{0.33\textwidth}
        \begin{block}{SMOTE}
            Genera muestras sintéticas interpolando.
            \[
            \mathbf{x}_{\text{nuevo}} = \mathbf{x}_i + \lambda(\mathbf{x}_j - \mathbf{x}_i)
            \]
            Más robusto que oversampling simple.
        \end{block}
    \end{columns}
    % Basado en Pregunta 12
\end{frame}

\begin{frame}{Técnicas a Nivel de Algoritmo}
    \begin{block}{Pesos en la función de pérdida}
        \[
        J_{\text{pond}} = -\frac{1}{n} \sum \left[ w_1 y \log(\hat{p}) + w_0 (1-y) \log(1-\hat{p}) \right]
        \]
        Pesos típicamente inversamente proporcionales a la frecuencia.
    \end{block}
    \begin{block}{Ajuste del umbral}
        \begin{itemize}
            \item Disminuir umbral para favorecer recall (detectar más positivos).
            \item Aumentar umbral para favorecer precisión.
        \end{itemize}
    \end{block}
    % Basado en Pregunta 13 y 14
\end{frame}

\begin{frame}{Evaluación en Problemas Desbalanceados}
    \begin{itemize}
        \item Evitar accuracy.
        \item Usar matriz de confusión, precisión, recall, F1 (especialmente para la clase minoritaria).
        \item Curva Precision-Recall y su área (PR-AUC) más informativa que ROC-AUC.
    \end{itemize}
    % Basado en Pregunta 11
\end{frame}

% ============================================================================
% SECCIÓN 5: CASO DE ESTUDIO INTEGRADO
% ============================================================================
\section{Caso de Estudio: Diagnóstico Médico}

\begin{frame}{Contexto y Preprocesamiento}
    \begin{itemize}
        \item \textbf{Datos:} pacientes con características clínicas (edad, presión, glucosa...).
        \item \textbf{Objetivo:} detectar enfermedad rara (prevalencia 10\%).
        \item \textbf{Preprocesamiento:}
        \begin{itemize}
            \item Estandarización de variables numéricas.
            \item Imputación de nulos.
            \item División estratificada (70/30).
        \end{itemize}
    \end{itemize}
    % Basado en Pregunta 19
\end{frame}

\begin{frame}{Modelado y Evaluación}
    \begin{columns}[t]
        \column{0.5\textwidth}
        \textbf{Estrategias probadas:}
        \begin{enumerate}
            \item Sin tratamiento del desbalance.
            \item Con pesos de clase.
            \item Con SMOTE.
        \end{enumerate}
        \column{0.5\textwidth}
        \textbf{Métricas en prueba:}
        \begin{itemize}
            \item Matriz de confusión.
            \item Precisión, recall, F1 para clase positiva.
            \item AUC-ROC y AUC-PR.
        \end{itemize}
    \end{columns}
    \vspace{0.3cm}
    \begin{exampleblock}{Resultados esperados}
        Sin tratamiento: bajo recall. Con pesos o SMOTE: mejora recall (a costa de precisión). La elección depende del costo de FN vs FP.
    \end{exampleblock}
\end{frame}

% ============================================================================
% SECCIÓN 6: CASO AVANZADO: DETECCIÓN DE FRAUDE (Pregunta 20)
% ============================================================================
\section{Caso Avanzado: Detección de Fraude}

\begin{frame}{Diseño de un sistema antifraude}
    \begin{itemize}
        \item 10M transacciones/día, fraude 0.1\%.
        \item Requisitos: interpretabilidad (regulatorio), velocidad (ms), buena detección.
    \end{itemize}
    \begin{block}{Feature engineering específico}
        \begin{enumerate}
            \item Log del monto.
            \item Hora del día (codificación cíclica seno/coseno).
            \item Agregadas por usuario: frecuencia últimas 24h, monto promedio, desviación.
            \item Ratios: monto actual vs promedio del usuario.
            \item Puntuación de anomalía (Isolation Forest) como predictor.
        \end{enumerate}
    \end{block}
    % Basado en Pregunta 20
\end{frame}

\begin{frame}{Manejo del desbalance extremo}
    \begin{itemize}
        \item \textbf{Nivel datos:} SMOTE (cuidado con ruido) o undersampling.
        \item \textbf{Nivel algoritmo:} pesos de clase y ajuste de umbral.
        \item \textbf{Métricas:} recall, precisión, F1, curva PR.
    \end{itemize}
    \vspace{0.3cm}
    \begin{block}{Monitoreo en producción}
        \begin{itemize}
            \item Técnicas: recall $\leftrightarrow$ fraude no detectado; precisión $\leftrightarrow$ transacciones legítimas bloqueadas.
            \item Impacto de falsos positivos: cliente insatisfecho, costo de atención.
        \end{itemize}
    \end{block}
\end{frame}

% ============================================================================
% SECCIÓN 7: RESUMEN Y CONEXIÓN
% ============================================================================
\section{Resumen y Conexión}

\begin{frame}{Lo que hemos aprendido}
    \begin{exampleblock}{Conceptos clave}
        \begin{itemize}
            \item Regresión logística: modelo probabilístico, sigmoide, log-loss, límite lineal.
            \item Interpretación mediante odds ratios.
            \item Evaluación: matriz de confusión, precisión, recall, F1, ROC/AUC.
            \item Manejo de clases desbalanceadas: muestreo (SMOTE), pesos, umbral.
            \item Aplicación a casos reales (spam, fraude, diagnóstico médico).
        \end{itemize}
    \end{exampleblock}

\end{frame}

% --- Diapositiva final ---
\begin{frame}
    \centering
    \Huge \textbf{¡Gracias!}

\end{frame}

\end{document}